\section{Introduction}

Text detoxification rewrites harmful or insulting text into a non-toxic form while
preserving meaning and language. This is already difficult in a single language. It is
harder across languages with different scripts, morphological systems, tokenization
behavior, data availability, and culturally specific forms of harmful speech. A model
may encode a toxic-to-detox transformation in a way that is separable in one language
but less accessible in another.

Mechanistic interpretability provides tools for studying this problem. Linear probes
test whether hidden states distinguish toxic inputs from detoxified targets. Sparse
autoencoders (SAEs) identify candidate latent features associated with the
transformation. Activation steering tests whether modifying a representation changes
generated text. These tools answer different questions. A high probe score establishes
decodability, not causality. A steering-induced output change establishes an
intervention effect, not necessarily improved non-toxicity or utility. A multilingual
study must keep these claims separate.

We investigate three research questions:

\noindent\textbf{RQ1: Circuit structure.} Do multilingual models use shared or
language-specific internal features when transforming toxic text into detoxified text?

\noindent\textbf{RQ2: Cross-lingual transfer.} Do directions identified in
higher-resource languages transfer to Yoruba and isiXhosa, and how does this compare
with within-language directions?

\noindent\textbf{RQ3: Causal steering.} Can activation steering improve
detoxification while preserving semantic content, fluency, and output diversity?

Our contributions are as follows:
\begin{enumerate}[leftmargin=*]
    \item We provide multi-seed evidence that toxic-input and detoxified-target
    representations are linearly decodable in Aya-23-8B and Qwen2.5-7B-Instruct.
    \item We show a consistent reduction in transfer from seven higher-resource
    languages to isiXhosa and Yoruba under the stored probe protocol.
    \item We trace layerwise transfer behavior in Qwen and analyze SAE candidate
    features and recovered steering sweeps across model families.
    \item We separate decodability, intervention-induced generation change, and
    validated detoxification quality through an explicit evaluation framework.
\end{enumerate}
